\section{Nonlinear selectors preserving an actual pure square}
\label{np19-section}

Imposing the different-owner congruences on the root parameter before
squaring preserves the exact power representation. The following
ordinary proof concerns exponent two and fixed marked data. It does
not concern the original moving prime-index block.

Retain \(\mathcal O=\mathbb Z[\zeta]\) and
\(\zeta^2-\zeta+1=0\). For \(k\ge1\) integral put
\[
 (3k+\zeta)^2=A_k+C_k\zeta,\qquad c_k=A_k+C_k,\qquad
 T(k)=A_kC_kc_k.
\]

\begin{proposition}[The actual square and its full boundary]
\label{np19-actual}
The three arms are positive and pairwise coprime, with
\[
 A_k=9k^2-1,\qquad C_k=6k+1,\qquad c_k=3k(3k+2),
\]
and
\begin{equation}
 T(k)=(3k-1)(3k+1)(6k+1)(3k)(3k+2).
 \label{np19-boundary}
\end{equation}
The input norm \(9k^2+3k+1\) is \(1\pmod3\). Every prime divisor
of the full boundary \(T(k)\) is a unit for this input norm.
\end{proposition}
\begin{proof}
Squaring in the stated ring gives the coordinate formulas and
factorization. If a positive integer divides both \(A_k,C_k\),
it divides \(4A_k-C_k^2=-2C_k-3\) and \(2C_k\), hence divides
\(3\). Since \(A_k\equiv-1\pmod3\), their gcd is one.
The other arm gcds follow from \(c_k=A_k+C_k\).

Norm multiplicativity gives
\[
 (9k^2+3k+1)^2=A_k^2+A_kC_k+C_k^2.
\]
At a prime dividing any one of the three pairwise coprime arms,
the right side is the square of a unit, including the case that
the sum arm vanishes. Thus the input norm is nonzero modulo that
prime. No factorization of the unmarked primes is assumed.
\end{proof}

\begin{theorem}[Exact finite packets with different actual owners]
\label{np19-packet}
Fix a finite nonempty set \(\mathcal S\) of primes \(q>5\).
For each \(q\), choose a positive integer \(k(q)\) with actual
depth \(e_q=v_q(T(k(q)))\ge1\). Choose a positive integer \(k_0\)
by CRT satisfying
\[
 k_0\equiv k(q)\pmod{q^{e_q+1}}\quad(q\in\mathcal S).
\]
Put
\[
 \mathcal E=\mathcal S\cup\{2,3,5\},\qquad
 F_p=v_p(T(k_0)),\qquad
 H=\prod_{p\in\mathcal E}p^{F_p+1},\quad
 K=\prod_{p\in\mathcal E}p^{F_p}.
\]
For every integer \(j\ge1\), set \(k_j=k_0+Hj\).
Then \((3k_j+\zeta)^2\) is a positive primitive unramified pure
square with exact depth \(F_p\) at every \(p\in\mathcal E\).
In particular its marked depths are \(e_q\), and its entire
selected signed cost is
\[
 \sum_{q\in\mathcal S}(e_q-3)\log q.
\]
\end{theorem}
\begin{proof}
The integer polynomial \(T\) preserves its value modulo each
\(q^{e_q+1}\). Fixing one digit beyond the actual valuation
therefore gives \(F_q=e_q\), and
\(\prod_{q\in\mathcal S}q^{e_q+1}\mid H\).
For every \(p\in\mathcal E\),
\(k_j\equiv k_0\pmod{p^{F_p+1}}\), so its exact depth remains
\(F_p\). Proposition~\ref{np19-actual} gives positivity,
primitivity and all norm-unit conclusions. The signed sum keeps
the terms of depths one and two as well as the positive ones.
\end{proof}

\begin{theorem}[Full squarefree sieve on the unchanged pure family]
\label{np19-density}
With all data in Theorem~\ref{np19-packet} fixed, there is a set
\(\mathcal G\) of positive integers of natural density
\[
 d(\mathcal G)=\prod_{p\notin\mathcal E}(1-5/p^2)\ge\frac16
\]
such that every \(j\in\mathcal G\) has
\[
 T(k_j)=KR_j,\qquad R_j\text{ squarefree},\qquad
 \gcd(K,R_j)=1,
\]
and all prime divisors of \(R_j\) lie outside \(\mathcal E\).
For the complete signed sum \(J(T)=\sum_{p\mid T}(v_p(T)-3)\log p\),
\begin{equation}
 J(T(k_j))=-2\log T(k_j)+3\log\frac K{\operatorname{rad}K}.
 \label{np19-signed}
\end{equation}
Along these good indices,
\[
 \frac{J(T(k_j))}{\log c_{k_j}}\longrightarrow-5,\qquad
 \frac{\log\operatorname{rad}(T(k_j))}{\log c_{k_j}}
 \longrightarrow\frac52.
\]
Every good index satisfies
\begin{equation}
 c_{k_j}\le
 \frac{15}{324^{2/5}}
 \left(\frac K{\operatorname{rad}K}\right)^{2/5}
 \operatorname{rad}(T(k_j))^{2/5}.
 \label{np19-radical}
\end{equation}
\end{theorem}
\begin{proof}
For every \(p\notin\mathcal E\), the five factors of
\eqref{np19-boundary} become affine polynomials in \(j\) with
\(p\)-unit slopes. In the coordinate \(a=3k\), their roots are
\[
 1,\quad -1,\quad -\tfrac12,\quad 0,\quad -2.
\]
They are distinct modulo \(p>3\): multiply them by two and note
that every nonzero difference between \(2,-2,-1,0,-4\) has
prime divisors only \(2,3\). The change \(a=3k_0+3Hj\) is
invertible modulo \(p^2\). Thus exactly five simple residue
classes modulo \(p^2\) are forbidden. Since the factors have
distinct roots already modulo \(p\), a square divisor \(p^2\)
of their product must divide one individual factor.

For \(1\le j\le X\), every positive factor is at most \(C_0X\),
where \(C_0=6(k_0+H)+1\). A prime square dividing a factor has
\(p\le\sqrt{C_0X}\). Counting its five residue classes gives the
complete large-prime estimate
\begin{equation}
 \begin{split}
 &\#\{1\le j\le X:\exists p>Y,\ p\notin\mathcal E,\
                    p^2\mid T(k_j)\}\\
 &\hspace{12mm}\le5X\sum_{p>Y}p^{-2}+5\sqrt{C_0X}.
 \end{split}
 \label{np19-tail}
\end{equation}
For a fixed finite cutoff \(Y\), CRT gives the finite product
density. Divide \eqref{np19-tail} by \(X\), first let \(X\) tend
to infinity, and then let \(Y\) tend to infinity. The upper and
lower finite-CRT bounds establish the exact natural density.
Since all omitted primes are at least seven,
\[
 \sum_{p\notin\mathcal E}\frac5{p^2}
 \le5\sum_{m\ge7}m^{-2}
 <5\int_6^\infty x^{-2}\,dx=\frac56.
\]
The finite inequality \(\prod(1-u_i)\ge1-\sum u_i\), followed
by the convergent product limit, proves the density lower bound.
This treats all prime squares; there is no unbounded-degree
squarefree-value assumption or unexamined large-prime range.

The frozen depths inside \(\mathcal E\) and squarefreeness
outside it give \(T=KR_j\) and
\[
 \operatorname{rad}T=T/(K/\operatorname{rad}K).
\]
This proves \eqref{np19-signed}, including every negative term.
The five affine factors give \(\log T=5\log j+O(1)\), while
\(\log c=2\log j+O(1)\), proving the two limits.

For all \(k\ge1\), the five lower bounds
\(2k,3k,6k,3k,3k\) give \(T(k)\ge324k^5\), and
\(c_k\le15k^2\). Eliminating \(k\) and using the radical identity
gives \eqref{np19-radical}.
\end{proof}

The set of resulting root indices has density \(d(\mathcal G)/H>0\)
among all positive integers. It consists of actual pure squares
satisfying the fixed marked data, not merely compatible local labels.

\paragraph{A positive-cost packet and the boundary of the conclusion.}
Take the two actual owners
\[
 \begin{split}
 k(19)&=4973918=21720+2\cdot19^5,\\
 k(13)&=4888691=61882+13^6.
 \end{split}
\]
Their second arms are \(19^4\cdot229\) and \(13^5\cdot79\).
At \(q>3\) dividing \(6k+1\), the other coordinate is
\(-3/4\pmod q\), and the input norm is \(3/4\pmod q\).
The whole depths are therefore exactly four and five.
The other owner has no hit: \(k(19)\equiv1\pmod{13}\) and
\(k(13)\equiv10\pmod{19}\), and substitution in the five
factors proves all are units at that other prime. The selected
cost is \(\log19+2\log13>0\). Theorem~\ref{np19-density}
then supplies infinitely many actual pure-square outputs with
this fixed positive packet and the full negative credit in
\eqref{np19-signed}. The earlier finite selector certificate
checks the displayed owner arithmetic; it is not an infinite
squarefree-density proof.

The residual coefficient remains one and the exponent is two,
so the earlier convention \(\lambda=1/n\) gives the fixed value
\(1/2\), which does not approach zero. All packet primes, depths,
\(k_0,H,K,C_0\) are fixed in the density limit. The size of
\(K/\operatorname{rad}K\), the progression step and the effective
tail threshold can depend on the full packet. The construction
does not give a moving-packet estimate in the original
\(n>324\) prime, \(B=n^4\) block.

The elementary arithmetic of these literal squares is separately
amenable to the scoped selector Lean module. The CRT, complete
sieve, natural density and asymptotic statements here are ordinary
proofs, not claims of complete formalization. The original far
signed tail, its exceptions, dependent-root complement, general
residuals and full ABC coverage remain open. The linear-factor
proof above cannot be applied to an arbitrary higher-degree
boundary without first checking its actual factors.
